\scanpage{201}\mspage{150}
Considérons le carré cartésien

\diagram{\begin{tikzcd}[ampersand replacement=\&,row sep=1.2em,column sep=2em]
U_{/t}\arrow[r,hook,"i_{/t}"]\arrow[d,"p"'] \& X_{/t}\arrow[d,"q"]\\
U\arrow[r,hook,"i"'] \& X
\end{tikzcd}\(\,,\)}

nous allons prouver qu'il est $W$-cocartésien. Moyennant $W(7\mathrm{bis})$, il s'ensuivra que $i\in W$ implique $i_{/t}\in W$, et comme $X_{/t}$ est \inserted{$W$-}asphérique (ayant un objet final), cela signifie que $U_{/t}$ est $W$-asphérique, ce qu'on veut prouver.

Il suffit, en vertu du cor. 2, de prouver que le carré est cocartésien dans $\mathrm{Cat}$, puisque $i_{/t}$ est une imm. ouverte ($i$ l'étant), et $X_{/t}\to X$ est un iso. local donc \inserted{$W$-}lisse. On voit que $X_{/t}\setminus U_{/t}\simeq e$ (catégorie \deleted{\unsure{finie|finale}} \inserted{ponctuelle}), \deleted{\unsure{donc $X_{/t}$ déduit de $U_{/t}$ en}} \inserted{d'ailleurs l'unique \unsure{objet|élément} de $X_{/t}\setminus U_{/t}$} est l'élém. final $t_0:t\xrightarrow{\mathrm{id}_t}t$ de $X_{/t}$. Ainsi $H_{t_0}\simeq e_{U_{/t}}$ (préfaisceau final), et $p_!(H_{t_0})\simeq \unsure{\mathbb H|H}'_t\overset{\mathrm{déf}}{=}(u\longmapsto\mathrm{Hom}_X(u,t))$, \deleted{\unsure{donc}} ce qui caractérise la somme amalgamée dans $\mathrm{Cat}$\ldots.

Le cor. 4 n'est autre que le corollaire non démontré de la prop \unsure{7} (p. \unsure{85}).

Voici une variante du cor. 4

\uline{Coro. 5} \inserted{(Moyennant $W(7\mathrm{bis})^{*}$)}
% slanted
\margin{forme forte.} Soit $i:U\hookrightarrow X$ une imm. ouverte, soit $t_0\in\mathrm{Ob}\,T$ un point \uline{fermé} de $T=X\setminus U$, tel que $\{t_0\}$
% archive sheet number: 201
